% 03_theory.tex -- Section 3: theoretical properties of the SMC-BVS posterior.
% Owner: Writing Department, Task Lead 3.  Draft of 2026-11-27.
% Source: Theory TH-SONG-ASSUMPTIONS package v1.3 (origin/dept/theory @ 99b11e4f,
%   Theory_work/Dept/song_assumptions/THEOREM_DEFINITIONS.md and ASSUMPTIONS_AND_PROOF_PATHS.md),
%   with the Stage-1 senior generality recommendations (council/SENIOR_FABLE.md, SENIOR_ASTRA.md):
%   population-only selection margin, no proof constants in the assumptions, primitive form of the
%   omitted-signal assumption.  Statements will be refined when Theory's final pass lands.
% Map to Theory labels:  A1 A-GL, A2 A-SPARSE, A3 A-DIM, A4 A-SUBG, A5 A-NEC, A6 A-BS,
%   A7 A-MARGIN-U (senior population form), A8 A-MOM, A9 A-MIS(iii) primitive form.
%   Theorems: TD (thm:dist) -> thm:th-dist; TS-b -> thm:th-select; TM-a0 / TM-a,TM-b -> thm:th-mis (a)/(b).
% Owner scope (HQ 2026-11-27/28): Song-Liang counterparts only -- true distribution, true model,
%   misspecification.  Pruned: sure screening (TS-a, A-MARGIN-H), the deletion converse, the beta-min
%   proposition (kept as a sentence of justification), the almost-sure corollary TD-as.

\section{Theoretical properties}\label{sec:theory}

This section gives the large-sample properties of the posterior of Section~\ref{sec:model}. There are three. The
posterior recovers the true density of the data at a nearly parametric rate (Theorem~\ref{thm:th-dist}); it puts
almost all its mass on the true model, so that both the highest posterior model and the median probability model
select it (Theorem~\ref{thm:th-select}); and when some true predictors are not measured, it targets the best linear
approximation built from the predictors that are (Theorem~\ref{thm:th-mis}). Each result has a counterpart in
\citet{SongLiang2015}, listed in Table~\ref{tab:sl-map}. The results concern the exact posterior, which the sampler of
Section~\ref{sec:model} approximates; they allow $p$ to exceed $n$ by any fixed power.

\subsection{Setting and notation}\label{sec:theory-setting}

The observations $(x_i,y_i)$, $i=1,\dots,n$, are independent copies of a pair $(x,y)$ with $x\in\R^p$. The
number of predictors $p=p_n$ may grow with $n$, and so may the law of $(x,y)$. A constant is \emph{fixed} if it does
not depend on $n$. The true model is
\begin{equation}\label{eq:th-truth}
  y=\alpha^*+x^\top\bstar+\varepsilon,\qquad \varepsilon\sim\Normal(0,\sigs^2),
\end{equation}
with $\Sz=\supp\bstar$ and $\sz=\abs{\Sz}$. We write $\nu$ for the law of $x$, $\Sig$ for its covariance matrix, and
$\zeta_j=(x_j-\E x_j)/\sd(x_j)$ for the standardised predictors. The true joint density of $(x,y)$, with respect to
$\nu(dx)\,dy$, is $f^*(x,y)=\phi(y;\alpha^*+x^\top\bstar,\sigs)$, where $\phi(\cdot\,;\mu,\sigma)$ is the
$\Normal(\mu,\sigma^2)$ density.

A posterior draw $(S,\beta,\sigma^2)$, mapped back to the original units, defines the conditional law
$\Normal(\bar y+(x-\bar x)^\top\beta,\sigma^2)$ of $y$ given $x$. Let $f$ be the joint density of $(x,y)$ when $x\sim\nu$
and $y$ given $x$ has this law. Let $\Pi(\cdot\mid\Dn)$ be the posterior of $(S,\beta,\sigma^2)$, and
$\P^*$ the law of the data. As in \citet{SongLiang2015}, we measure distance by the Hellinger distance,
\[
  d_H(f,f^*)=\Big(1-\int\!\!\int\sqrt{f(x,y)\,f^*(x,y)}\;\nu(dx)\,dy\Big)^{1/2},
\]
which is their distance divided by $\sqrt2$.

Two quantities drive the selection results. The first is the \emph{residual signal} of a model $S$,
\begin{equation}\label{eq:th-Delta}
  \Delta_n(S)=n\,\inf_{a\in\R,\,b\in\R^{\abs S}}\E\big(x^\top\bstar-a-x_S^\top b\big)^2 ,
\end{equation}
which is $n$ times the error of the best linear prediction of the true mean from the predictors in $S$. It is zero when
$S\supseteq\Sz$. The second is the \emph{price} that the prior and the marginal likelihood charge the true model for
$d$ predictors that a smaller model leaves out,
\begin{equation}\label{eq:th-price}
  \mathrm{pen}_n(d)=\log\frac{\pi(\Sz\setminus D)}{\pi(\Sz)}+\frac d2\log(n+1)
  =\log\frac{\binom{2p-1}{\sz}}{\binom{2p-1}{\sz-d}}+\frac d2\log(n+1),\qquad \abs D=d.
\end{equation}
The first term is the prior odds, the second the Occam factor of the $g$-prior with $g=n$. Near the true model each
left-out predictor costs about $\log(2p/\sz)+\tfrac12\log n$, so $\mathrm{pen}_n(d)$ grows like $d\log p$.

\subsection{Consistency of the fitted density}\label{sec:theory-density}

We first state the conditions on the data.

\begin{assumption}[Gaussian linear model]\label{a:th-gl}
The pairs $(x_i,y_i)$ are independent copies of $(x,y)$, and $y$ follows~\eqref{eq:th-truth} with $\varepsilon$
independent of $x$ and $\sigs^2>0$. The quantities $\alpha^*$, $\bstar$, $\sigs$ and $\nu$ may depend on $n$.
\end{assumption}

Assumption~\ref{a:th-gl} is model~(1) of \citet{SongLiang2015} with an intercept and a random design. It makes the
marginal likelihood of every model exact, so that the posterior can be studied model by model. It is the standard
working model for a continuous response; it fails for counts and for errors whose variance changes with~$x$.
Section~\ref{sec:theory-mis} removes it in two ways.

\begin{assumption}[Sparsity]\label{a:th-sparse}
$\sz\le s$ for a fixed integer $s\ge0$.
\end{assumption}

Assumption~\ref{a:th-sparse} is the strict sparsity condition~(4) of \citet{SongLiang2015}; \citet{NarisettyHe2014}
also fix the size of the true model. It says that a few predictors matter however large $p$ becomes. It is what makes
the number of competing models of the relevant size polynomial in $p$, and so controllable.

\begin{assumption}[Dimension]\label{a:th-dim}
$2\le p\le n^A$ for a fixed $A\ge1$.
\end{assumption}

Assumption~\ref{a:th-dim} lets $p$ exceed $n$ by any fixed power. \citet{JohnsonRossell2012} prove selection
consistency for $p\le n$, and \citet{ShinBhattacharyaJohnson2018} note that their results continue to hold when $p$
grows polynomially. \citet{SongLiang2015} allow
faster growth, $p<\exp(Cn^\alpha)$ with $\alpha<1$, at the cost of a slower rate. Under Assumption~\ref{a:th-dim},
$\log p$ and $\log n$ are of the same order, so a single rate $\rn$ describes the posterior. Whether our results hold for
$\log p=o(\sqrt n)$ is open.

In the next two assumptions, a \emph{sparse set} is a set of at most $\min(p,k_n)$ predictors, for a fixed sequence
$k_n$ with $k_n/\log n\to\infty$.

\begin{assumption}[Sub-Gaussian predictors]\label{a:th-subg}
Every predictor has positive variance, and for a fixed $L\ge1$,
$\E\exp(\lambda u^\top\zeta)\le\exp\{L^2\lambda^2\Var(u^\top\zeta)/2\}$ for every $\lambda\in\R$ and every $u\in\R^p$
supported on a sparse set.
\end{assumption}

Assumption~\ref{a:th-subg} asks that every combination of a few predictors have tails no heavier than a Gaussian with
the same variance. Condition~1 of \citet{SongLiang2015} plays the same role: it bounds exponential moments of averages
of up to $\bar r_n$ predictors. Ours is scaled by the variance of each combination, as in \citet{RudelsonZhou2013}. It
holds for Gaussian rows, and for rows $x=m+\Lambda\xi$ with $\xi$ of independent, bounded, standardised coordinates. Its
use is that the sample correlations of all sparse sets are then close to the population ones at the same time, so
population conditions carry over to the data.

\begin{assumption}[No near-exact collinearity]\label{a:th-nec}
$n\,\lmin\{\Corr(x_T)\}\ge c_\psi$ for every sparse set $T$, for a fixed $c_\psi>0$.
\end{assumption}

Assumption~\ref{a:th-nec} allows the smallest sparse eigenvalue to tend to zero, but no faster than $1/n$. It excludes
only predictors that are nearly exact linear combinations of a few others, such as duplicated columns. The scale $1/n$
is set by the prior, not by the data: the slab $\Normal(0,n\sigma^2(X_S^\top X_S+I)^{-1})$ adds a unit ridge to a Gram
matrix of order $n\Corr(x_S)$, and the Occam factor $(n+1)^{-\abs S/2}$ of a model is intact as long as the eigenvalues
of that matrix are not small against the ridge. \citet{SongLiang2015} recommend a ridge of the same order
($c_n\propto n^{-1}$), and their identifiability condition, in its Gaussian form, is a floor on the smallest eigenvalue
that also tends to zero. Fixed floors on sparse eigenvalues
\citep{BickelRitovTsybakov2009,NarisettyHe2014} are a special case. No upper bound on the eigenvalues is needed, so
strong common factors are allowed.

\begin{assumption}[Bounded effects]\label{a:th-bs}
$\abs{\beta^*_j}\,\sd(x_j)\le\bar b\,\sigs$ for every $j$, for a fixed $\bar b>0$.
\end{assumption}

Assumption~\ref{a:th-bs} says that no single standardised effect is unboundedly large compared with the noise. It is
implied by the approximate sparsity condition~(3) of \citet{SongLiang2015}, $\sum_j\abs{\beta^*_j}<\infty$ for
standardised predictors, when the noise level is fixed; \citet{YangWainwrightJordan2016} need a bound of the same type,
in terms of $g$, for the $g$-prior. Its use
is technical: it keeps the noise variance of the fitted models in a fixed range and the ridge bias negligible. It
places no lower bound on any effect.

We now state the first result. Let $\epsilon_n=\sqrt{M\log n/n}$ for a level $M\ge1$.

\begin{theorem}[Consistency of the fitted density]\label{thm:th-dist}
Assume \ref{a:th-gl}--\ref{a:th-bs}. There are fixed constants $c_H>0$ and $M_0$, depending only on
$(s,A,L,c_\psi,\bar b)$, such that for every $M\ge M_0$ and all sufficiently large $n$,
\[
  \P^*\big[\,\Pi\{d_H(f,f^*)>\epsilon_n\mid\Dn\}>n^{-c_HM}\,\big]\le n^{-c_HM}.
\]
\end{theorem}
% Draft status (visible line below to be removed at submission).
\noindent\textit{[Draft status: proved and independently checked (delivered thm:dist; registry row 2.1).]}

Theorem~\ref{thm:th-dist} has the form of Theorems~1 and~2 of \citet{SongLiang2015}: their bound
$\exp(-c_1n\epsilon_n^2)$ equals $n^{-c_1M}$ at our rate. Three differences matter. First, no lower bound on the
effects is needed: effects too small to be selected do not move the fitted density by more than $\epsilon_n$. Second,
nothing is tuned. \citet{SongLiang2015} require conditions (9)--(14) on the prior inclusion probability $r_n/P_n$,
the size cap $\bar r_n$ and the prior covariance $V_\xi$. Our cap $K=\min(p,\lfloor n/2\rfloor)$ is far larger than
their (10) allows, but, as they remark (p.~951), a cap can be dropped when the prior probability of large models decays
fast enough; under the beta-binomial prior it decays like $2^{-k}$. Third, the rate is explicit.

\subsection{Selection of the true model}\label{sec:theory-selection}

Density consistency does not imply selection: a true predictor can be replaced by a combination of other predictors
with almost the same fit. \citet{SongLiang2015} rule this out with an identifiability condition (their Condition~3).
We use the residual signal~\eqref{eq:th-Delta}, and ask it to beat the price~\eqref{eq:th-price}.

A rival of $\Sz$ is a model $S=(\Sz\setminus D)\cup N$ that leaves out a
nonempty set $D\subseteq\Sz$ of $d$ true predictors and adds a set $N$ of $m\ge0$ null predictors. Models that
contain $\Sz$ need no signal condition: they explain no more of the true mean than $\Sz$, and the prior and the Occam
factor charge them for every extra predictor. For a rival, the added nulls can absorb
part of the signal of the left-out predictors. We measure this by the \emph{leakage}
\begin{equation}\label{eq:th-leak}
  \rho_n(D,m)=1-\frac{\min_N\Delta_n\{(\Sz\setminus D)\cup N\}}{\Delta_n(\Sz\setminus D)},
\end{equation}
the largest share of the residual signal of $\Sz\setminus D$ that $m$ null predictors can explain; the minimum is
over sets $N$ of $m$ null predictors. The leakage is a population partial $R^2$, computable from $\Sig$ alone. It is
zero for $m=0$ and for independent predictors. Let $\omega(\rho)=\min\{1,(2\rho)^{1/2}\}$ and
$\bar m=\lceil(2A+1)s\rceil+1$.

\begin{assumption}[Selection margin]\label{a:th-margin}
$\sz\ge1$, and there is a fixed $\eta>0$ such that, for all sufficiently large $n$ and every nonempty $D\subseteq\Sz$,
\begin{align}
  \Delta_n(\Sz\setminus D)\ &\ge\ 2(1+\eta)^2\,\sigs^2\,\mathrm{pen}_n(\abs D), \label{eq:th-margin-del}\\
  \Big[\frac{\min_N\Delta_n\{(\Sz\setminus D)\cup N\}}{\sigs^2}\Big]^{1/2}
  \ &\ge\ (1+\eta)\Big[\big\{2\,\mathrm{pen}_n(\abs D)\big\}^{1/2}
      +\omega\{\rho_n(D,m)\}\Big\{2\log\binom{p-\sz}{m}\Big\}^{1/2}\Big] \label{eq:th-margin-sub}
\end{align}
for every $m$ with $1\le m\le\min(\bar m,p-\sz)$ and $\sz-\abs D+m\le K$.
\end{assumption}

Condition~\eqref{eq:th-margin-del} is the deletion margin: the signal lost by leaving out $D$ must exceed twice the
price of $D$. The constant~$2$ is the $\tfrac12$ of the Gaussian log-likelihood. The log posterior odds of $\Sz$
against $\Sz\setminus D$ are about $\hat\Delta/(2\sigs^2)-\mathrm{pen}_n(\abs D)$, where $\hat\Delta$ is the sample
version of $\Delta_n(\Sz\setminus D)$, so the break-even point is $\Delta_n=2\sigs^2\mathrm{pen}_n$. This is also why
the price has the form of the extended BIC penalty \citep[Section~2.1]{SongLiang2015}. Condition~\eqref{eq:th-margin-sub}
is the substitution margin. Its last term is the most that the best of $\binom{p-\sz}{m}$ null sets can gain by
chance. The factor $\omega\in[0,1]$ says which fraction of the worst case applies: $\omega=0$ when the nulls cannot
imitate the missing signal, $\omega=1$ when nothing is assumed about them, and then the term is about
$\sqrt{2m\log p}$. Rivals with more than $\bar m$ nulls need no condition. Such a rival fits no better than
$\Sz\cup N$, and the Occam factor of the added nulls, of order $n^{-m/2}$ up to logarithmic factors, outweighs the
price of any deletion, which is at most about $s(A+\tfrac12)\log n$.

Assumption~\ref{a:th-margin} is a condition on the law of $(x,y)$ only; it involves no tuning and no sample quantity.
Conditions of this type are standard. Condition~3 of \citet{SongLiang2015} asks that no combination of other
predictors mimic a true one; ours says how well it may mimic it. \citet{NarisettyHe2014} (Condition~4.4) ask the
residual signal of every model that misses a true predictor to exceed $5(1+\delta)\sigs^2\abs{\Sz}\log(\sqrt n\vee p)$,
and \citet{LuoChen2013} ask its ratio to $\log p$ to diverge; ours charges each missed predictor its own price, with
constant~$2$, and allows that ratio to stay bounded. The order $\log p$ cannot be avoided: below a threshold of this
order no method recovers the support \citep{Wainwright2009,FletcherRanganGoyal2009}.

The familiar beta-min form is a special case. Let $\varphi^2\le\lmin\{\Corr(x_T)\}$ for every set $T$ of at most
$s+\bar m$ predictors, and $b_{\min}=\min_{j\in\Sz}\abs{\beta^*_j}\sd(x_j)/\sigs$. Since
$\Delta_n(S)\ge n\varphi^2\sum_{j\in D}\beta_j^{*2}\Var(x_j)$ for every rival,
$\mathrm{pen}_n(d)\le d\log(2p\sqrt{n+1})$, $\omega\le1$ and $\log\binom{p-\sz}{m}\le m\log p$,
Assumption~\ref{a:th-margin} holds whenever
\[
  b_{\min}\ \ge\ \frac{(1+\eta)(1+\sqrt{\bar m})}{\varphi}\,\Big\{\frac{2\log\big(2p\sqrt{n+1}\big)}{n}\Big\}^{1/2}.
\]
For independent predictors the leakage is zero and the requirement drops to
$\beta_j^{*2}\Var(x_j)\ge2(1+\eta)^2\sigs^2\log(2p\sqrt{n+1})/n$ for every $j\in\Sz$: each standardised effect must
exceed the noise level by the factor $\{2\log(2p\sqrt n)/n\}^{1/2}$. Assumption~\ref{a:th-margin} itself is sharper,
because it charges correlation only through the leakage.

\begin{theorem}[Selection consistency]\label{thm:th-select}
Assume \ref{a:th-gl}--\ref{a:th-margin}. Then $\pis{\Sz}\to1$ in probability,
\[
  \max_{1\le j\le p}\big\lvert\PIP_j-\ind{j\in\Sz}\big\rvert\ \le\ 1-\pis{\Sz}\ \to\ 0\quad\text{in probability},
\]
and $\P(\HPM=\MPM=\Sz)\to1$.
\end{theorem}
\noindent\textit{[Draft status: target; proof path open (the superset sum and the mixed-rival odds remain to be
written). A delivered special case, with fixed effects and a fixed eigenvalue floor, is proved and independently
checked (thm:post-random).]}

Theorem~\ref{thm:th-select} is the counterpart of Theorem~S in the supplement of \citet{SongLiang2015}, which shows
that the posterior odds of every rival against the true model tend to zero, and of their Theorem~3, which shows that
the inclusion probability of every identifiable true predictor tends to one. It gives more than both: the posterior
mass of the true model tends to one, every inclusion probability converges to $0$ or $1$, and the two usual point
estimates, the highest posterior model and the median probability model of \citet{BarbieriBerger2004}, both equal
$\Sz$. Their eigenvalue condition~(22), which they note is hardly satisfied, and its random-design replacement~(26)
are not needed; Assumptions~\ref{a:th-subg} and~\ref{a:th-nec} take their place.

\subsection{Misspecified models}\label{sec:theory-mis}

Suppose the method is run on the response and a set $C\subset\{1,\dots,p\}$ of candidate predictors that may miss some
true ones. As in Section~2.3 of \citet{SongLiang2015}, the set $C$ is fixed before the data are seen, or chosen on
independent data. The best the method can do is to find the best linear predictor of $y$ from $x_C$. Let
\[
  (\alpha_C,\beta^C)=\arg\min_{a,b}\E(y-a-x_C^\top b)^2,\qquad
  \sigma_C^2=\E(y-\alpha_C-x_C^\top\beta^C)^2,\qquad \Sz^C=\supp\beta^C,
\]
and let $f_0(x_C,y)=\phi(y;\alpha_C+x_C^\top\beta^C,\sigma_C)$. This $f_0$ is the Kullback--Leibler projection of the
truth on the Gaussian linear models in $x_C$, display~(17) of \citet{SongLiang2015}, and $\beta^C$ is their~(19). The
set $\Sz^C$ may contain null predictors that stand in for missing ones, and may lose true predictors whose effect is
cancelled.

Two conditions cover two kinds of misspecification. The first allows any law of $(x_C,y)$ in fixed dimension. Let
$Z=(1,x_C^\top)^\top$ and $u=y-\alpha_C-x_C^\top\beta^C$.

\begin{assumption}[Fixed candidate dimension]\label{a:th-mom}
The law of $(x,y)$ does not depend on $n$; $\abs C\ge2$ is fixed; $\E y^2<\infty$, $\E\norm{x_C}^2<\infty$,
$\Var(x_C)$ is nonsingular and $\sigma_C^2>0$.
\end{assumption}

Assumption~\ref{a:th-mom} is the setting of Theorem~4 of \citet{SongLiang2015}: a fixed subset of predictors with
finite second moments (their Condition~4). It asks nothing of the form of $\E(y\mid x_C)$ or of the error law: the
response may be non-Gaussian, heteroscedastic or nonlinear in $x_C$. Second moments are what is needed for the best
linear predictor to exist, so the assumption is close to minimal.

In high dimension, misspecification of the error law can make Bayesian linear regression overfit
\citep{GrunwaldVanOmmen2017}. We therefore allow only the misspecification that comes from leaving out predictors. Let
$O=\Sz\setminus C$ be the true predictors that are left out.

\begin{assumption}[Omitted predictors]\label{a:th-mis}
The left-out signal satisfies $x_O^\top\bstar_O=a+x_C^\top\gamma+e$, where $e\sim\Normal(0,\tau^2)$ is independent of
$x_C$ and $\varepsilon$, and $\gamma$ has at most $s$ nonzero entries. Assumptions~\ref{a:th-dim}--\ref{a:th-bs} hold for
the candidate problem, with $(x_C,\beta^C,\sigma_C^2)$ in place of $(x,\bstar,\sigs^2)$.
\end{assumption}

Under Assumptions~\ref{a:th-gl} and~\ref{a:th-mis}, $\beta^C=\bstar_C+\gamma$ and $\sigma_C^2=\sigs^2+\tau^2$ (the
minimiser with at most $2s$ nonzero entries, unique by Assumption~\ref{a:th-nec}): given $x_C$ the data are
again Gaussian and linear, and the left-out signal acts as extra noise and as extra effects on the candidates that
predict it. The assumption holds whenever $(x_C,x_O)$ is jointly Gaussian; then $\gamma$ is supported on the candidates
adjacent to $O$ in the conditional-independence graph of the predictors, so it is sparse when each left-out predictor
has few neighbours among the candidates. It fails under strong common factors, where the projection spreads over many
candidates. The sparsity of $\gamma$ keeps the pseudo-true model $\Sz^C$, of size at most $2s$, within the reach of
Theorems~\ref{thm:th-dist} and~\ref{thm:th-select}.

\begin{theorem}[Consistency under misspecification]\label{thm:th-mis}
Let $\Pi(\cdot\mid\Dn^C)$ be the posterior computed from $(y,X_C)$, and $\epsilon_n$ as in
Theorem~\ref{thm:th-dist}.
\begin{enumerate}[label=(\alph*)]
  \item \emph{Any response law, fixed dimension.} Under Assumption~\ref{a:th-mom}, for every $\epsilon>0$,
    $\Pi\{d_H(f,f_0)>\epsilon\mid\Dn^C\}\to0$ and $\Pi(S\supseteq\Sz^C\mid\Dn^C)\to1$ almost surely. If in addition
    $\E(Z_j^2u^2)<\infty$ for every coordinate $Z_j$ of $Z$, then $\pi(\Sz^C\mid\Dn^C)\to1$ almost surely.
  \item \emph{Omitted predictors, high dimension.} Under Assumptions~\ref{a:th-gl} and~\ref{a:th-mis}, there are fixed $c>0$ and $M_0$ such that, for every
    $M\ge M_0$ and all sufficiently large $n$,
    \[
      \P^*\big[\,\Pi\{d_H(f,f_0)>\epsilon_n\mid\Dn^C\}>n^{-cM}\,\big]\le n^{-cM}.
    \]
    If moreover $\Sz^C\ne\emptyset$ and Assumption~\ref{a:th-margin} holds for the candidate problem, then
    $\pi(\Sz^C\mid\Dn^C)\to1$ in probability and $\P(\HPM=\MPM=\Sz^C)\to1$.
\end{enumerate}
\end{theorem}
\noindent\textit{[Draft status: (a) target, low risk (direct law-of-large-numbers and iterated-logarithm argument over
finitely many models); (b) target: density part low--medium risk, selection part inherits the open steps of
Theorem~\ref{thm:th-select}. A delivered analogue of the selection part, for jointly Gaussian predictors with
detectable projected effects, is proved and independently checked (thm:omitted).]}

Part~(a) is the counterpart of Theorem~4 of \citet{SongLiang2015}, which gives the density statement. The inclusion
statement is their Corollary~2. The last statement is stronger: the posterior selects exactly the support of the best
linear predictor, so a candidate is kept if and only if $\beta^C_j\neq0$, which their discussion of Corollary~2
anticipates but their display does not give. Part~(b) brings the rates and the exact selection of
Sections~\ref{sec:theory-density}--\ref{sec:theory-selection} to high-dimensional candidate sets. In both parts, a
null predictor that stands in for a missing true one is selected, and the noise variance absorbs what the candidates
cannot explain. These are the properties a method needs when, as is usual in practice, not every relevant predictor
has been measured.

The scope of the three results is deliberate. They take the number of true predictors fixed, the noise Gaussian
with constant variance (except in Theorem~\ref{thm:th-mis}(a)), and $p$ polynomial in $n$. They describe the exact
posterior, not a finite run of the sampler. Within this scope they hold for any correlation structure that meets
Assumptions~\ref{a:th-subg}--\ref{a:th-margin}, common factors included, and they require nothing to be tuned.

\begin{table}[ht]
\centering\small
\caption{Results of Section~2 of \citet{SongLiang2015} and their counterparts here.}\label{tab:sl-map}
\begin{tabular}{@{}p{0.30\linewidth}p{0.27\linewidth}p{0.35\linewidth}@{}}
\toprule
\citet{SongLiang2015} & Here & Difference \\
\midrule
Theorems 1--2, Corollary 1 (density) & Theorem~\ref{thm:th-dist} & explicit rate; no tuning sequences; $p\le n^A$ \\
Theorem 3, Theorem S (inclusion, MAP model) & Theorem~\ref{thm:th-select} & posterior mass of $\Sz$; HPM and MPM; all PIPs \\
Theorem 4 (misspecified subset) & Theorem~\ref{thm:th-mis}(a) & also exact support of $\beta^C$ \\
Corollary 2 (pseudo-true support) & Theorem~\ref{thm:th-mis}(a)--(b) & inclusion and exclusion; high-dimensional $C$ \\
\bottomrule
\end{tabular}
\end{table}
